% GENERATED FILE. Edit canonical lessons, then run npm run generate:books.
% canonical-source-hash: fnv1a64:dc9382d91dd86cf3
% canonical-lessons: FA-C372-tasir, FA-C372-tasir-gozashtan, FA-C372-nashi-shodan, FA-C372-mojeb-shodan, FA-C372-vakonesh

\chapter{Effects and reactions}
\label{ch:fa-effects-and-reactions}
\hlchaptermodality{372}

\begin{chapteropening}
\textbf{By the end of this chapter:} \emph{I can say what influenced me, what something stems from, what it gave rise to and how people reacted.}

Complete the last lesson of chapter 372: I can say what influenced me, what something stems from, what it gave rise to and how people reacted.
\end{chapteropening}

\section[\texorpdfstring{\fa{تاثیر} (ta'sir)}{تاثیر (ta'sir)}]{\fa{تاثیر} (ta'sir) — an influence, an effect}

\label{lesson:FA-C372-tasir}



\begin{quote}
Before the new one: say the Persian for therefore, then the Persian for considering.
\end{quote}

\subsection*{You'll want to know: \fa{تاثیر}}
\textbf{\fa{تاثیر}} (\emph{ta'sir}) — \textquotedblleft{}an influence, an effect\textquotedblright{}.

\begin{quote}
\fa{این} \fa{کتاب} \fa{تاثیر} \fa{زیادی} \fa{روی} \fa{من} \fa{داشت}.
\end{quote}

\begin{quote}
\emph{in ketâb ta'sir-e ziyâdi ruye man dâsht.}
\end{quote}

\begin{quote}
This book had a big influence on me.
\end{quote}

\emph{Ta'sir dâshtan bar} or \emph{ruye} — to have an effect on. Plural \emph{ta'sirât}. Many writers put a hamza on its alef; here it is written plain.

\begin{cousinweb}
\textbf{\fa{تاثیر}} is usually traced to Arabic \emph{ta'thīr} \textquotedblleft{}effect\textquotedblright{}, a verbal noun usually explained as leaving a trace; Persian says the Arabic \emph{th} as \emph{s}.
\end{cousinweb}

\begin{grammarlens}[title={What you've built}]
One more word for telling, explaining and describing at B1.
\end{grammarlens}

\subsection*{Your turn}
\begin{itemize}
\raggedright
  \item \emph{Say it:} \emph{ta'sir}
  \item \emph{Say it:} \emph{ta'sir}, once more
  \item \emph{Say it:} a sentence of your own with \emph{ta'sir}
  \item \emph{Recall:} say the Persian for therefore, then the Persian for considering, then say \emph{ta'sir} again
\end{itemize}

\begin{tcolorbox}[breakable,colback=teal!4,colframe=teal!35!black,title={Before you move on}]
What does \fa{تاثیر} mean? (\textquotedblleft{}An influence, an effect\textquotedblright{}.) Say it once more.
\end{tcolorbox}
\section[\texorpdfstring{\fa{تاثیر} \fa{گذاشتن} (ta'sir gozâshtan)}{تاثیر گذاشتن (ta'sir gozâshtan)}]{\fa{تاثیر} \fa{گذاشتن} (ta'sir gozâshtan) — to influence, to have an impact}

\label{lesson:FA-C372-tasir-gozashtan}



\begin{quote}
Before the new one: say the Persian for considering, then the Persian for an influence.
\end{quote}

\subsection*{You'll want to know: \fa{تاثیر} \fa{گذاشتن}}
\textbf{\fa{تاثیر} \fa{گذاشتن}} (\emph{ta'sir gozâshtan}) — \textquotedblleft{}to influence, to have an impact\textquotedblright{}.

\begin{quote}
\fa{معلمم} \fa{روی} \fa{زندگی} \fa{من} \fa{تاثیر} \fa{گذاشت}.
\end{quote}

\begin{quote}
\emph{mo'allemam ruye zendegi-ye man ta'sir gozâsht.}
\end{quote}

\begin{quote}
My teacher had an impact on my life.
\end{quote}

\emph{Ta'sir} + \emph{gozâshtan} \textquotedblleft{}to put\textquotedblright{} (present stem \emph{gozâr-}). Takes \emph{ruye} or \emph{bar}.

\begin{cousinweb}
\textbf{\fa{تاثیر} \fa{گذاشتن}} is literally \textquotedblleft{}to put an effect\textquotedblright{}: \emph{ta'sir} is usually traced to Arabic \emph{ta'thīr}, and \emph{gozâshtan} \textquotedblleft{}to put, to leave\textquotedblright{} is a native Persian verb.
\end{cousinweb}

\begin{grammarlens}[title={What you've built}]
One more word for telling, explaining and describing at B1.
\end{grammarlens}

\subsection*{Your turn}
\begin{itemize}
\raggedright
  \item \emph{Say it:} \emph{ta'sir gozâshtan}
  \item \emph{Say it:} \emph{ta'sir gozâshtan}, once more
  \item \emph{Say it:} a sentence of your own with \emph{ta'sir gozâshtan}
  \item \emph{Recall:} say the Persian for considering, then the Persian for an influence, then say \emph{ta'sir gozâshtan} again
\end{itemize}

\begin{tcolorbox}[breakable,colback=teal!4,colframe=teal!35!black,title={Before you move on}]
What does \fa{تاثیر} \fa{گذاشتن} mean? (\textquotedblleft{}To influence, to have an impact\textquotedblright{}.) Say it once more.
\end{tcolorbox}
\section[\texorpdfstring{\fa{ناشی} \fa{شدن} (nâshi shodan)}{ناشی شدن (nâshi shodan)}]{\fa{ناشی} \fa{شدن} (nâshi shodan) — to stem from, to arise from}

\label{lesson:FA-C372-nashi-shodan}



\begin{quote}
Before the new one: say the Persian for an influence, then the Persian for to influence.
\end{quote}

\subsection*{You'll want to know: \fa{ناشی} \fa{شدن}}
\textbf{\fa{ناشی} \fa{شدن}} (\emph{nâshi shodan}) — \textquotedblleft{}to stem from, to arise from\textquotedblright{}.

\begin{quote}
\fa{سردرد} \fa{او} \fa{از} \fa{خستگی} \fa{ناشی} \fa{میشود}.
\end{quote}

\begin{quote}
\emph{sardard-e u az khastegi nâshi mishavad.}
\end{quote}

\begin{quote}
His headache comes from tiredness.
\end{quote}

Takes \emph{az}: \emph{az chizi nâshi shodan}. Formal.

\begin{cousinweb}
\textbf{\fa{ناشی}} is usually traced to Arabic \emph{nāshi'} \textquotedblleft{}arising, growing out of\textquotedblright{}; \emph{shodan} \textquotedblleft{}to become\textquotedblright{} is the inherited Persian verb.
\end{cousinweb}

\begin{grammarlens}[title={What you've built}]
One more word for telling, explaining and describing at B1.
\end{grammarlens}

\subsection*{Your turn}
\begin{itemize}
\raggedright
  \item \emph{Say it:} \emph{nâshi shodan}
  \item \emph{Say it:} \emph{nâshi shodan}, once more
  \item \emph{Say it:} a sentence of your own with \emph{nâshi shodan}
  \item \emph{Recall:} say the Persian for an influence, then the Persian for to influence, then say \emph{nâshi shodan} again
\end{itemize}

\begin{tcolorbox}[breakable,colback=teal!4,colframe=teal!35!black,title={Before you move on}]
What does \fa{ناشی} \fa{شدن} mean? (\textquotedblleft{}To stem from, to arise from\textquotedblright{}.) Say it once more.
\end{tcolorbox}
\section[\texorpdfstring{\fa{موجب} \fa{شدن} (mojeb shodan)}{موجب شدن (mojeb shodan)}]{\fa{موجب} \fa{شدن} (mojeb shodan) — to give rise to, to cause}

\label{lesson:FA-C372-mojeb-shodan}



\begin{quote}
Before the new one: say the Persian for to influence, then the Persian for to stem from.
\end{quote}

\subsection*{You'll want to know: \fa{موجب} \fa{شدن}}
\textbf{\fa{موجب} \fa{شدن}} (\emph{mojeb shodan}) — \textquotedblleft{}to give rise to, to cause\textquotedblright{}.

\begin{quote}
\fa{کار} \fa{زیاد} \fa{موجب} \fa{شد} \fa{که} \fa{او} \fa{مریض} \fa{شود}.
\end{quote}

\begin{quote}
\emph{kâr-e ziyâd mojeb shod ke u mariz shavad.}
\end{quote}

\begin{quote}
Too much work made him ill.
\end{quote}

\emph{Mojeb} \textquotedblleft{}cause\textquotedblright{} + \emph{shodan}; formal, like \emph{bâ'es shodan}. \emph{Ke} + subjunctive follows.

\begin{cousinweb}
\textbf{\fa{موجب}} is usually traced to Arabic \emph{mūjib} \textquotedblleft{}that which makes something necessary\textquotedblright{}; with \emph{shodan} \textquotedblleft{}to become\textquotedblright{} it means to be the cause.
\end{cousinweb}

\begin{grammarlens}[title={What you've built}]
One more word for telling, explaining and describing at B1.
\end{grammarlens}

\subsection*{Your turn}
\begin{itemize}
\raggedright
  \item \emph{Say it:} \emph{mojeb shodan}
  \item \emph{Say it:} \emph{mojeb shodan}, once more
  \item \emph{Say it:} a sentence of your own with \emph{mojeb shodan}
  \item \emph{Recall:} say the Persian for to influence, then the Persian for to stem from, then say \emph{mojeb shodan} again
\end{itemize}

\begin{tcolorbox}[breakable,colback=teal!4,colframe=teal!35!black,title={Before you move on}]
What does \fa{موجب} \fa{شدن} mean? (\textquotedblleft{}To give rise to, to cause\textquotedblright{}.) Say it once more.
\end{tcolorbox}
\section[\texorpdfstring{\fa{واکنش} (vâkonesh)}{واکنش (vâkonesh)}]{\fa{واکنش} (vâkonesh) — a reaction}

\label{lesson:FA-C372-vakonesh}



\begin{quote}
Before the new one: say the Persian for to stem from, then the Persian for to give rise to.
\end{quote}

\subsection*{You'll want to know: \fa{واکنش}}
\textbf{\fa{واکنش}} (\emph{vâkonesh}) — \textquotedblleft{}a reaction\textquotedblright{}.

\begin{quote}
\fa{واکنش} \fa{مادرم} \fa{به} \fa{این} \fa{خبر} \fa{جالب} \fa{بود}.
\end{quote}

\begin{quote}
\emph{vâkonesh-e mâdaram be in khabar jâleb bud.}
\end{quote}

\begin{quote}
My mother's reaction to this news was interesting.
\end{quote}

A modern Persian coinage: \emph{vâ-} \textquotedblleft{}back\textquotedblright{} + \emph{konesh} \textquotedblleft{}action\textquotedblright{}. \emph{Vâkonesh neshân dâdan} — to react.

\begin{cousinweb}
\textbf{\fa{واکنش}} is a modern native coinage built from \emph{vâ-} \textquotedblleft{}back, again\textquotedblright{} + \emph{konesh} \textquotedblleft{}action\textquotedblright{}, built on \emph{kon-}, the present stem of \emph{kardan} \textquotedblleft{}to do\textquotedblright{}. A reaction is an action back.
\end{cousinweb}

\begin{grammarlens}[title={What you've built}]
One more word for telling, explaining and describing at B1.
\end{grammarlens}

\subsection*{Your turn}
\begin{itemize}
\raggedright
  \item \emph{Say it:} \emph{vâkonesh}
  \item \emph{Say it:} \emph{vâkonesh}, once more
  \item \emph{Say it:} a sentence of your own with \emph{vâkonesh}
  \item \emph{Recall:} say the Persian for to stem from, then the Persian for to give rise to, then say \emph{vâkonesh} again
\end{itemize}

\begin{tcolorbox}[breakable,colback=teal!4,colframe=teal!35!black,title={Before you move on}]
What does \fa{واکنش} mean? (\textquotedblleft{}A reaction\textquotedblright{}.) Say it once more.
\end{tcolorbox}
